% ECONOMICS_PROBLEM: {"id": "L52-201", "title": "Industrial-policy spillovers", "jel_code": "L52", "source_id": "OP-0346"}
% !TeX program = xelatex
\documentclass[11pt,a4paper]{article}
\usepackage[margin=23mm]{geometry}
\usepackage{fontspec}
\setmainfont{Latin Modern Roman}
\usepackage{amsmath,amssymb,mathtools}
\usepackage{xcolor,enumitem,booktabs,longtable,array,xurl,hyperref}
\definecolor{muted}{HTML}{53636C}
\hypersetup{colorlinks=true,urlcolor=blue,linkcolor=blue}
\setlength{\parindent}{0pt}
\setlength{\parskip}{5pt}
\newcommand{\R}{\mathbb R}
\newcommand{\E}{\mathbb E}
\newcommand{\N}{\mathbb N}
\newcommand{\ind}{\mathbf 1}
\newcommand{\Var}{\operatorname{Var}}
\newcommand{\Cov}{\operatorname{Cov}}
\newcommand{\supp}{\operatorname{supp}}
\DeclareMathOperator*{\argmin}{arg\,min}
\DeclareMathOperator*{\argmax}{arg\,max}
\newcommand{\entrymeta}[3]{{\sffamily\small\color{muted}\textbf{\texttt{#1}}\quad|\quad #2\par #3\par}}
\newcommand{\Problemref}[1]{\texttt{#1}}
\newcommand{\JELref}[2]{\texttt{#2} (JEL \texttt{#1})}
\begin{document}
\section*{Industrial-policy spillovers}
\textbf{Problem ID:} \texttt{L52-201}\quad \textbf{JEL:} \texttt{L52}\par
% BEGIN ECONOMICS STATEMENT
\label{problem:OP-0346}
{\sffamily\small\color{muted}Original subject group: Trade, globalization and geoeconomics\par}
\label{definitions:problem:30007628}
\textbf{Audit classification:} Published research agenda. The full published review discusses unresolved mechanisms and economy-wide counterfactual/spillover evaluation. The catalogue international subsidy/retaliation model is editorial; strategic subsidy races are a game-theoretic variant.
\subsection*{Definitions and precise research target}
Fix countries, industries and a subsidy policy \(a\) targeted at production, investment or learning. Firms' technologies permit specified learning, knowledge spillovers and market-power mechanisms. Let \(Z_{cjt}(a)\) be productivity and \(W_c(a)\) discounted household welfare in country \(c\), including fiscal financing and equilibrium foreign retaliation. Define policy effects relative to an otherwise specified path \(0\):
\[\Delta Z_{cj}(h)=E[Z_{cj,t+h}(do(a))-Z_{cj,t+h}(do(0))],\qquad \Delta W_c=W_c(a)-W_c(0).\]
The task is to identify conditions yielding durable domestic productivity gains and characterize foreign spillovers, rents and retaliation costs. Separate true efficiency gains from output expansion caused by transfers or market-share displacement. State the shock or assignment design that disentangles policy from anticipated industry growth, and determine which spillovers remain unobserved. Include endogenous subsidy races rather than treating foreign policy as permanently fixed when evaluating global welfare. Estimate or bound learning externalities and financing costs with independently disciplined data. Report transportability across institutional capacity and market structure. A positive treated-firm effect does not establish a national welfare gain, and a national gain does not establish positive aggregate international welfare.

\textbf{Additional formal definitions and scope (editorial).} Specify subsidy instruments in money per output unit, ad valorem rates or investment grants, including eligibility and fiscal funding. Productivity \(Z\) is physical output per defined input bundle or an estimated production-function residual, not revenue productivity unless explicitly chosen. A learning law such as \(Z_{cj,t+1}=\Phi_{cj}(Z_t,q_t,a_t,\varepsilon_{t+1})\) must distinguish private learning from knowledge externalities. Foreign spillovers are \(\Delta W_d\) and \(\Delta Z_{dj}\) for untargeted countries \(d\), allowing trade-price, knowledge and policy-response channels. Global welfare fixes country weights and avoids treating rents paid by one agent to another as net resource creation. If retaliation is endogenous, specify its decision problem and equilibrium selection; this strategic extension falls within game-theoretic economic policy analysis. Holding foreign policies fixed gives a different unilateral causal target.

For the identification part, declare an admissible class \(\Theta\) of structural models, including preferences, technologies, shock laws, measurement/sampling rules and any equilibrium selector. If \(O\) denotes the complete observable history and \(T(\theta)\) the target defined above, the identified set is \(\mathcal I_T=\{T(\theta):\theta\in\Theta,\ \mathcal L_\theta(O)=\mathcal L_{\rm obs}(O)\}\); point identification means this set is a singleton. A \(do(a)\) comparison replaces stated policy/technology equations while fixing the declared exogenous law and re-solving the remaining equations. These definitions do not assert that an identification theorem holds without further restrictions.
\subsection*{Variants and assumption changes}
\begin{enumerate}
\item \textbf{Unilateral policy, fixed foreign rules (editorial).} Identify learning or productivity effects and welfare at home and abroad under a prespecified subsidy assignment.
\item \textbf{Policy interaction and spillovers (editorial; strategic extension).} Specify foreign reaction rules or a policy game and compare global welfare under common weights; do not label market-share transfers as aggregate productivity gains.
\end{enumerate}
\subsection*{References and source audit}
\begin{enumerate}
\item Publisher and full PDF verify Juhasz, Lane and Rodrik, 2024, Annual Review of Economics 16:213--242 and DOI. Locator: Sections 4.4--4.5, printed pp.228--231; especially economy-wide counterfactual discussion, p.231. Support: Mechanisms and economy-wide spillovers agenda. Full published PDF accessible via publisher doi/pdf endpoint. URL: \url{https://www.annualreviews.org/doi/pdf/10.1146/annurev-economics-081023-024638}.
\end{enumerate}
\begin{enumerate}\raggedright
\item\hypertarget{definitions:source:30007628:1}{} Réka Juhász; Nathan Lane; Dani Rodrik. 2024. \emph{The New Economics of Industrial Policy}. Sections 4.4--4.5, printed pp.228--231; especially economy-wide counterfactual discussion, p.231.
\url{https://www.annualreviews.org/doi/pdf/10.1146/annurev-economics-081023-024638}.
DOI: \url{https://doi.org/10.1146/annurev-economics-081023-024638}.
\end{enumerate}

\subsection*{Common conventions: Source mathematical conventions}

These conventions apply to each entry unless explicitly replaced.

\textbf{Models and identification.} A primitive model \(m\) specifies all probability laws, feasible choices, information, equilibrium and measurement rules used by its target. The admissible class \(\mathcal M\) is fixed independently of outcomes. If \(P_O(m)\) is its observable probability law and \(\tau(m)\) the target, its identified set at observable law \(P\) is
\[
 \mathcal I_\tau(P)=\{\tau(m):m\in\mathcal M,\ P_O(m)=P\}.
\]
Point identification means that this set is a singleton. An empty set signals incompatibility of data and assumptions. Coordinate bounds need not describe the joint set. Bounds are sharp only if every retained target is attainable or arbitrarily approachable by compatible models. Finite bounds require explicit restrictions on unobserved quantities; unrestricted confounding, hidden wealth, extrapolation or arbitrary equilibrium selection can yield uninformative sets.

\textbf{Causal interventions.} A potential outcome \(Y_i(p)\) is the outcome under a complete policy regime \(p\), including timing, financing, information and market spillovers. The target population and units are fixed before treatment. With interference, use the complete assignment vector or a justified exposure mapping; individual no-interference notation alone cannot identify an economy-wide intervention. A difference of mean potential outcomes does not require the cross-world joint distribution. Conditioning on post-treatment migration, survival, enrollment or employment changes the estimand and needs separate justification.

\textbf{Statistical statements.} An estimator, confidence set, forecast or test specifies a sampling law, model class and loss/coverage criterion. Uniform \((1-\alpha)\)-coverage means \(\inf_{m\in\mathcal M}\Pr_m\{\tau(m)\in C_n\}\geq1-\alpha\), or an explicitly asymptotic version. The whole parameter space trivially has coverage, so informativeness requires a stated width, risk or power objective. A forecasting improvement is relative to a named benchmark, information set and evaluation population.

\textbf{Equilibrium and optimization.} An equilibrium comprises feasible plans, optimality, beliefs where needed, budget constraints and all market/resource clearing. The primitives or a stated benchmark define each correspondence \(\mathcal E(p;m)\). Nonemptiness is assumed or proved, never inferred just from compact policy sets. A selected-equilibrium target uses a declared selection rule; a robust target quantifies over all permitted equilibria. Use suprema or infima unless attainment is established. An \(\varepsilon\)-optimal policy is feasible and within \(\varepsilon\) of the finite optimum value. Report an empty feasible set separately.

\textbf{Welfare and accounting.} Normative utility, interpersonal normalization, social weights, numeraire, discounting and the use of public receipts are primitives. Behavioral utility need not coincide with the welfare criterion. Household welfare includes ownership income and rents once; adding profits again to compensating variation generally double-counts them. A monetary equivalent is defined by an explicit transfer and price convention, with monotonicity and range conditions. Average effects, mechanism contributions, transfers and welfare losses are different targets. Infinite sums require convergence and dynamic feasible plans require terminal or no-Ponzi conditions.

\textbf{Source versions.} A working-paper date, revision date and journal publication year are distinguished. A DOI for a journal version is not automatically the DOI of an earlier manuscript. Locators refer to the stated version and printed pages where specified. A metadata-only check does not verify a claim about a full-text conclusion. Every entry's source subsection narrows the provenance claim accordingly.
% END ECONOMICS STATEMENT
\end{document}
